% PLAN1-2026-11-11 additional independent/optional results excluded from 46-page core
% These are source excerpts, not a compilable manuscript.
% SOURCE BLOB: d7bb4f7d3c51d1a050d9f082e03726f342543621

% [SECTION 8] Reduction and two-letter results; may later retain only semilattice reduction in #2
\paragraph{Reduction to semilattices.}
Let $\mathbf{SL}$ be the pseudovariety of finite semilattices (commutative
idempotent semigroups).  For a nonempty finite alphabet $\Sigma$, the
\emph{content typing} $c_\Sigma$ maps $w\in\Sigma^+$ to the set of letters
occurring in $w$, a morphism onto the nonempty subsets of $\Sigma$ under union.

\begin{lemma}[Coarser typings inherit characteristic samples]
\label{lem:typing-monotone}
Let $s\preceq t$ be typings, $f,r\ge1$, and $L\in\Cmcf{f}{s}$.  For every
finite $K\subseteq L$,
\[
 L\bigl(\mathcal B^{\le r}_{f,t}(K)\bigr)\subseteq
 L\bigl(\mathcal B^{\le r}_{f,s}(K)\bigr)\subseteq L,
\]
and likewise for $\mathcal B_{f,t}$ and $\mathcal B_{f,s}$.  Hence every
characteristic sample of $L$ for $\mathcal B^{\le r}_{f,t}$ (resp.\
$\mathcal B_{f,t}$) is one for $\mathcal B^{\le r}_{f,s}$ (resp.\
$\mathcal B_{f,s}$).
\end{lemma}

\begin{proof}
Both grammars have the same nonterminals and the same Start, Const and Comp
rules (Definition~\ref{def:reconstruction}).  A Link rule of the $t$-grammar
joins observed tuples with a common context and equal $t$-types; since
$\ker t\subseteq\ker s$, their $s$-types are equal, so it is a Link rule of the
$s$-grammar.  This gives the first inclusion; the second is
Proposition~\ref{prop:reconstruction-soundness}.  If $F$ is characteristic for
the $t$-operator and $F\subseteq K\subseteq L$, the left-hand side equals $L$,
hence so does the middle term.
\end{proof}

\begin{proposition}[Reduction of the fan-out boundary to semilattices]
\label{prop:sl-reduction}
Fix $f\ge1$ and $r_{\max}\ge1$, and read uniform polynomial characteristic-data
length at fan-out $f$ and rank $r_{\max}$
(Definition~\ref{def:uniform-polydata}).
\begin{enumerate}[label=\textup{(\roman*)},leftmargin=*]
\item Every pseudovariety $\mathbf V\not\subseteq\mathbf{LG}$ contains
$\mathbf{SL}$.  Consequently, the pseudovarieties with uniform polynomial
characteristic-data length are exactly those contained in $\mathbf{LG}$ if and
only if $\mathbf{SL}$ does not have it.
\item $\mathbf{SL}$ has uniform polynomial characteristic-data length if and
only if, for every nonempty finite alphabet $\Sigma$, the content typing
$c_\Sigma$ satisfies the condition of Definition~\ref{def:uniform-polydata}.
\end{enumerate}
\end{proposition}

\begin{proof}
(i) A finite semigroup $\mathsf S\notin\mathbf{LG}$ is the positive image of
its evaluation typing over the alphabet $\mathsf S$, so it has a subsemigroup
isomorphic to $U_1$ (Theorem~\ref{thm:lg-characterization}); hence
$U_1\in\mathbf V$.  Every finite semilattice $\mathsf H$ embeds into
$U_1^{\mathsf H}$ via $x\mapsto(\chi_z(x))_{z\in\mathsf H}$, where $\chi_z(x)=1$ if
$zx=z$ and $\chi_z(x)=\zeta$ otherwise (a morphism, since $zxy=z$ iff $zx=z$ and
$zy=z$), so $\mathbf{SL}\subseteq\mathbf V$.  The property of
Definition~\ref{def:uniform-polydata} passes from a pseudovariety to its
subpseudovarieties, because fewer typings are quantified.  If $\mathbf{SL}$
lacks it, so does every $\mathbf V\not\subseteq\mathbf{LG}$, while every
$\mathbf V\subseteq\mathbf{LG}$ has it by Theorem~\ref{thm:local-group-polydata}.
Conversely, $U_1\notin\mathbf{LG}$, so $\mathbf{SL}\not\subseteq\mathbf{LG}$.

(ii) The image of $c_\Sigma$ is a semilattice, which gives one direction.  For
the other, let $s:\Sigma^+\twoheadrightarrow\mathsf S$ with $\mathsf S\in\mathbf{SL}$.
Since $\mathsf S$ is commutative and idempotent, $s(w)$ is the product of
$s(a)$ over the letters $a$ occurring in $w$, so $s\preceq c_\Sigma$ and
$\Cmcf{f}{s}\subseteq\Cmcf{f}{c_\Sigma}$ (item~(a) after
Definition~\ref{def:class-cmcf}).  For $L\in\Cmcf{f}{s}$ and a presentation $G$,
a characteristic sample for $\mathcal B^{\le r_{\max}}_{f,c_\Sigma}$ of length at
most $p_{c_\Sigma}(\|G\|,\bar\theta_G)$ is one for
$\mathcal B^{\le r_{\max}}_{f,s}$ by Lemma~\ref{lem:typing-monotone}; take
$p_{s}:=p_{c_\Sigma}$.
\end{proof}

By Theorem~\ref{thm:fanout-one-data}, $\mathbf{SL}$ lacks the property at
fan-out one when $r_{\max}\ge2$, although over two letters every typing in
$\mathbf{SL}$ has polynomial data there (Corollary~\ref{cor:two-letter-fanout-one}).
At fan-out two the question is open (Open Problem~\ref{op:fanout-two}), and the
next theorem settles two-letter alphabets at every fan-out $f\ge2$.  Rare-type lower bounds transfer to content
typings: for $s\preceq c_\Sigma$ every content type contained in an $s$-type $m$
has shortest word at least $\mu_m(L)$, computed for $s$.

\paragraph{Two letters.}
Call $L$ \emph{Parikh-saturated} if membership of a word of $\Sigma^+$ in $L$
depends only on its Parikh vector.

\begin{theorem}[Content typings on two letters at every fan-out $f\ge2$]
\label{thm:two-letter-content}
Let $\Sigma=\{\mathtt a,\mathtt b\}$, $c:=c_\Sigma$, $f\ge2$, $r_{\max}\ge2$, and put
$D_f:=(6f+1)3^f$ and
$C_{f,r_{\max}}:=(r_{\max}+1)\bigl(4r_{\max}(3f+1)+D_f\bigr)+3$, and let
$H_m:=\sum_{k=1}^m1/k\le1+\ln m$ \textup($H_0:=0$\textup).  For every
$L\in\Cmcf{f}{c}$ and every reduced good fan-out-$\le f$, rank-$\le r_{\max}$
presentation $G$ of $L\cap\Sigma^+$,
$\tau(G,c)\le C_{f,r_{\max}}|V|^2\bar\theta_G^{\,4}H_{\bar\theta_G}$, and the operators
$\mathcal B^{\le r_{\max}}_{f,c}$ and $\mathcal B_{f,c}$ have a common
characteristic sample $K$ with
\[
 \|K\|_+\le4^{fr_{\max}}\bigl(3+r_{\max}C_{f,r_{\max}}\bigr)\|G\|^3\bar\theta_G^{\,4}H_{\bar\theta_G}
 +\mathbf 1_{\{\eps\in L\}}.
\]
If $L$ is not Parikh-saturated, then
$t_{\mathrm{typ}}(G,c)\le D_f(r_{\max}+1)|V|\bar\theta_G$ and one may take
$\|K\|_+\le4^{fr_{\max}}\bigl(3+r_{\max}D_f(r_{\max}+1)\bigr)\|G\|^2\bar\theta_G
+\mathbf 1_{\{\eps\in L\}}$.
\end{theorem}

\begin{corollary}[Semilattice typings on two letters]
\label{cor:two-letter-semilattice}
Let $f\ge2$, $r_{\max}\ge2$ and let $s:\{\mathtt a,\mathtt b\}^+\twoheadrightarrow\mathsf S$ be a
typing with $\mathsf S\in\mathbf{SL}$, for instance $s_{U,2}$
\textup($\mathtt a\mapsto1$, $\mathtt b\mapsto\zeta$\textup).  Then the bound
required of each typing in Definition~\ref{def:uniform-polydata} holds for this
fixed typing at fan-out $f$ and rank $r_{\max}$, with the bound of
Theorem~\ref{thm:two-letter-content}, for
$\mathcal B^{\le r_{\max}}_{f,s}$ and for $\mathcal B_{f,s}$.
\end{corollary}

\begin{proof}
As in the proof of Proposition~\ref{prop:sl-reduction}(ii), $s\preceq c$ and
$\Cmcf{f}{s}\subseteq\Cmcf{f}{c}$; apply Theorem~\ref{thm:two-letter-content} and
Lemma~\ref{lem:typing-monotone}.
\end{proof}

The proof is in Appendix~\ref{app:two-letter-content}.  It uses the
conditions of arities one and two, which hold at every fan-out $f\ge2$, through
runs of one letter: once $x^4$ occurs in a word of
$L$ for a letter $x$, membership depends only on the $x$-skeleton (each maximal
run of $x$ replaced by one $x$) and the number of letters $x$, and two runs of the other letter next to a run $xx$ either
force Parikh saturation or merge into a forbidden longer run.  Without
saturation this bounds the antichains of heavy nodes in a shortest typed
derivation; with saturation, node-minimal derivations of words with a
prescribed Parikh vector are analysed directly.  The theorem does not separate
fan-outs: by Corollary~\ref{cor:two-letter-fanout-one}, two-letter
semilattice typings have polynomial data at fan-out one as well, so over two
letters they have polynomial data at every fan-out.  For three or more
letters a long component of content $\{\mathtt a,\mathtt b\}$ need not contain a
long run, and the question is open.

% [CONCLUSION] zeta capacity and finite foreign-letter shortening, independent of A-C
Rare types need many foreign letters.  For $M\subseteq\Sigma$ write
$\#_{\overline M}(w)$ for the number of occurrences in $w$ of letters outside $M$, and
$\kappa_{\overline M}(L):=\sup\{\#_{\overline M}(w):w\in L\}$.  For the content typing,
$\mu_M(L)$ is the length of a shortest word of $L$ with content exactly $M$; for a
typing $s$ onto $U_1=\{1,\zeta\}$ the words of type $1$ are those of $\Sigma_1^+$,
$\Sigma_1:=\{a\in\Sigma:s(a)=1\}$.

\begin{proposition}[Short words from few foreign letters]
\label{prop:zeta-capacity}
Let $G$ be a reduced $\eps$-free nondeleting MCFG of any fan-out and rank
$\le r_{\max}$, let $L:=L(G)$, let $\emptyset\ne M\subseteq\Sigma$, and suppose
$\kappa:=\kappa_{\overline M}(L)<\infty$.
\begin{enumerate}[label=\textup{(\roman*)},leftmargin=*]
\item If $L$ has a word with content $M$, then
$\mu_M(L)\le(2|M|-1+\kappa)\,\|G\|\bigl(\|G\|+r_{\max}\bar\theta_G\bigr)$.
\item If $L\cap M^+\ne\emptyset$, its shortest word has length at most
$\max\{\bar\theta_G,\ \kappa\,\|G\|(\|G\|+r_{\max}\bar\theta_G)\}$.  In particular,
for a typing onto $U_1$, $\mu_1(L)$ obeys this bound with $M=\Sigma_1$.
\end{enumerate}
\end{proposition}

\begin{proof}
For a nonterminal $A$ let $\lambda(A)$ be the least total length of a tuple of
$L_A(G)$ over $M$ ($\infty$ if there is none) and
$\operatorname{cap}(A):=\sup\{\#_{\overline M}(\tuple y):\tuple y\in L_A(G)\}$.  Since
$G$ is reduced, each $A$ has an outside context in $L$, so
$\operatorname{cap}(A)\le\kappa$; since $G$ is nondeleting, a derivation keeps every
letter of its subtrees, so $\operatorname{cap}(A)\ge\sum_j\operatorname{cap}(B_j)$ for
every rule $A\to\rho(B_1,\ldots,B_m)$.

(i) Among the derivation trees of words of $L$ of length $\mu_M(L)$ and content
$M$ choose one, $T$, with fewest nodes, and mark one occurrence of each letter of
$M$ in its word.  Call a change of $T$ \emph{admissible} if it keeps the marks and
adds only letters of $M$; the new word then has content $M$, so it is not
shorter, and if it is equally long it has at least as many nodes.  For a node
$\eta$ let $\mathrm{Mk}(\eta)$ be the set of marks written in its subtree.  The
nonempty sets $\mathrm{Mk}(\eta)$ form a laminar family on $|M|$ points, so they
take at most $2|M|-1$ values, and the nodes with one value form a chain.  If
two nodes of one chain had the same label, replacing the upper subtree by the
lower one would be admissible (the marks lie below the lower node, and the word
only loses letters) and would remove nodes.  So at most $(2|M|-1)|V|$ nodes have
marks below them.  A node $\eta$ without marks below derives a tuple
$\tuple x_\eta$ over $M$, and replacing its subtree by a derivation of a tuple of
length $\lambda(A_\eta)$ over $M$ is admissible; hence
$|\tuple x_\eta|=\lambda(A_\eta)$, $\lambda$ does not increase from such a node to
its children, and, as before, the labels on a path of such nodes are distinct.

Call an unmarked node \emph{heavy} if its label has $\lambda>\bar\theta_G$.  A rule
with left side $A$ is used in a successful derivation of a word of length at
most $\bar\theta_G$, and the tuple derived there from $A$ is shorter than
$\lambda(A)$, so it contains a letter outside $M$; hence heavy labels have
capacity at least $1$.  The unmarked parent of a heavy node is heavy, so the
heavy nodes form subtrees hanging from marked nodes.  If $q$ heavy nodes have no
heavy children, replacing their (disjoint) subtrees at once by derivations of
tuples with a letter outside $M$ gives a word of $L$ with at least $q$ such
letters; so $q\le\kappa$, and there are at most $\kappa|V|$ heavy nodes.  Every
terminal of the word is written by a marked or a heavy node, at most $\|G\|$ per
node, or lies below one of the at most $r_{\max}$ remaining children of such a
node, at most $\bar\theta_G$ per child.  With $|V|\le\|G\|$ this gives (i).

(ii) Argue in the same way without marks, starting from a tree of a shortest
word of $L\cap M^+$ with fewest nodes.  If the root is not heavy, the word has
length at most $\bar\theta_G$; otherwise all nodes are unmarked and the count
above applies.
\end{proof}

Neither substitutability nor the fan-out is used.  So a family that answers
Open Problem~\ref{op:fanout-two} positively through
Proposition~\ref{prop:rare-type}, for a content typing or for a typing onto
$U_1$, must have $\kappa_{\overline M}(L_n)$ superpolynomial in $\|G_n\|$ and
$\bar\theta_{G_n}$ for the rare content $M$.  The fan-out-one family $L_n$ of
Proposition~\ref{prop:three-letter-family} has $\kappa_{\overline{\{\mathtt a,\mathtt b\}}}(L_n)=2^n$,
so this condition does not by itself separate the fan-outs; at fan-out two the
known additional constraint is power transfer
(Lemma~\ref{lem:power-transfer}).

For a related quantity, the size of finite-monoid observations preserving
guarded tuple substitution, the arity at which it stabilizes over a
pseudovariety is $1$, $2$ or $\infty$ \cite{KuriyamaSCL2026}.

\paragraph{Tightness of the exponents.}
In $O(\|G\|^2\bar\theta_G)$, for window--group and local-group typings alike,
one factor $\|G\|$ counts the witness words and the other bounds root-to-leaf
paths in a $\lambda$-minimal derivation by the number of node labels
(Lemma~\ref{lem:heavy-node-count}).  The path factor is genuine for the witness
construction: the grammars $S\to SS\mid\mathtt a\mid\mathtt b\mid A_0\mid\cdots\mid A_n$,
$A_0\to\mathtt b$, $A_i\to\mathtt a\mid A_{i-1}\mathtt b$ ($1\le i\le n$) generate
$\{\mathtt a,\mathtt b\}^+$ with $\|G_n\|=10n+16$ and $\bar\theta_{G_n}=2$, but
$\mathtt b^{n+1}$ is the only yield of $A_n$ of $s_{1,1}$-type
$\langle\mathtt b,\mathtt b\rangle$, so $t_{\mathrm{typ}}(G_n,s_{1,1})\ge n+1$.  For this
target, however, $\tau(G_n,s_{1,1})\le2$
(Corollary~\ref{cor:distribution-representatives}), so the example bounds the
witness method, not the data.  For the trivial typing the upper bound
$O(\|G\|\bar\theta_G)$ is attained by the marker families
(Corollary~\ref{cor:trivial-exact}); for
every typing with $\mathsf S\in\mathbf{LG}$ the remaining gap is a single factor
$\|G\|$, a gap between the known examples and the bounds rather than a proved
ratio.  Neither lower-bound argument closes it: for a fixed typing the
rare-type argument yields at most $|\mathsf S|(t_{\mathrm{typ}}+1)$, and the marker
families consist of $O(\|G\|)$ words of length $\bar\theta_G=t_{\mathrm{typ}}$.
For noncommutative
groups the constant $M_{\mathbb G,f}$ is doubly exponential in $f$ because of the
Ramsey-type homogenization (Lemma~\ref{lem:order-homogenize}); $|\mathbb G|$ owners suffice
at $f=1$, and whether $\mathrm{poly}(|\mathbb G|^f)$ owners suffice for $f\ge2$, as for
abelian groups, is open.  At fan-out one, for isolated letters, the exponents
$(23,8)$ of Corollary~\ref{cor:two-letter-fanout-one} come from the counting
argument of Appendix~\ref{app:two-letter-fanout-one} and are far from the lower
bound $\bar\theta_G^{\,2}/4$ of Proposition~\ref{prop:binary-u1-quadratic}.
